\section{Literature review}

Authentication is one of the very important security services that can be applied in various networking domains .

The safe authentication mechanism utilizing blockchain technology was proposed by Yazdinejad et al. . Drones were designed to execute the planned deployment of the strategy in smart cities. At every stage of the process, this approach guaranteed the least amount of delays. A zone-based architecture was devised for a drone network, and a decentralized consensus mechanism tailored for remote drone use in smart cities was deployed.

Bera et al. introduced ACSUD-IoD, an innovative access control system designed to identify and thwart unwanted unmanned aerial vehicles (UAVs) within the IoD. The storing of transactional data within a private blockchain framework was enabled by the integration of a blockchain-based solution with ACSUD-IoD. This encompassed the delivery of secure, standardized data from an UAV to the ground station server. Consequently, the transactional data on the blockchain is verifiable. A formal security verification was performed utilizing the ``Automated Validation of Internet Security Protocols and Applications (AVISPA) tool'', alongside a comprehensive security evaluation. It illustrated that their method was adequately protected against several possible threats.

% Feng et al. proposed a ``cross-domain authentication protocol grounded in blockchain technology''. This system was designed to use 5G technology for diverse IoD applications. The aforementioned limits were duly acknowledged during the formulation of this plan with the aim of transcending them. Their methodology was based on a varied collection of signatures, all produced via threshold sharing. As a result, they successfully established a productive identity federation for collaborative domains.

Cho et al. developed an authentication mechanism for unmanned aerial vehicles (UAVs) to reduce security threats linked to unauthorized drones utilizing the IoD concept. Although their methodology reduced communication and computational requirements, yet their architecture was vulnerable to the ``Ephemeral Secret Leakage (ESL) attack under the CK-adversary model''. The method insufficiently protected the anonymity and untraceability of the participants. Another element that contributed to this issue was the absence of blockchain technology in their proposed strategy.

Gupta et al. presented a GaRuDa system, which might potentially denoted as the drone-based delivery system that operated on the blockchain technology. The integration of this system into the operations of Healthcare 5.0 applications was feasible. The IoT and blockchain technology were utilized in their approach to enable the swift and accurate distribution of medical supplies, which could be continuously monitored and recorded by many stakeholders. This was achieved by using a 5G-enabled Internet environment.

A pair of unique communication strategies for UAV environments were developed by Rodrigues et al. . Their scheme facilitated the establishment of a direct exchange of messages between two drones. Nevertheless, the main contractual arrangement has been altered within the framework of this strategy. In accordance with the CK-adversary concept, their scheme was not impervious to the possibility of an ESL attack. Moreover, their scheme lack support for the blockchain technology.

Ever  proposed an authentication system for IoT applications that used Elliptic Curve Cryptography (ECC). UAVs were considered to be mobile extensions of wireless sensor networks, operating within a hierarchical framework, according to their design. This particular design enabled the effective implementation of one-time user authentication for mobile sinks (UAVs), cluster chiefs, and sensor nodes. In contrast, their system was vulnerable to ``ESL attack under the CK-adversary model''. Moreover, their scheme did not ensure the maintenance and safeguarding of anonymity and untraceability. Another limitation of their scheme was absence of blockchain technology and it required more communication and computational costs.

Singh et al. examined the evolution and potential applications of the Internet of Drones. The advanced development of this technology has generated several apprehensions, among which the degree of security offered by autonomous robots has always been a prominent issue. Hence, they emphasized the most urgent security vulnerabilities and suggested that the most efficient approach to address these challenges would be to adopt state-of-the-art blockchain technology.

Xiong et al. introduced a secure collaborative computing system that implemented blockchain technology. They initially created a lightweight blockchain framework that was specifically designed for ``Unmanned Aerial Vehicle (UAV) Ad-Hoc Networks (UANET)''. Further, they introduced an improved ``Practical Byzantine Fault Tolerance (PBFT)'' consensus algorithm that was based on trust assessment.

Wang et al. introduced a mutual authentication method that was both simple and effective, and it exclusively relied on one-way hash algorithms and bitwise XOR operations. Additionally, the issue of a centralized trusted authority (TA) was mitigated by blockchain technology. The Real-or-Random model-based formal security analysis was employed. Further, an informal security proof was provided to prove the security of their proposed authentication mechanism. Further, Wang et al. introduced, ``BSIF: Blockchain-Based Secure, Interactive, and Fair Mobile Crowdsensing'' system. It was blockchain-based and was distinguished by its security, interactivity, and impartiality. These attributes were achieved through the integration of smart contracts and mobile devices. Yu et al. presented a ``Cross-domain Industrial IoT Based on Consortium Blockchain mechanism (CBDS) for the security of Industrial Internet of Things (IIoT). Further, they introduced consortium blockchain specifically to establish trust across IIoT domains.

Srinivas et al. developed an innovative authentication technique that was anonymous, lightweight, and relied on temporal credentials for Internet of Things (IoT)-based platforms. It was denoted as TCALAS. To enhance TCALAS, Ali et al. developed an improved version of TCALAS, referred to as iTCALAS for the secure communication of IoD.

Mishra et al. represented a framework for managing authentication and session keys using blockchain technology. This framework supported the integration of big data analytics capabilities for drones that operate on networks beyond 5G applications. Through a comprehensive security examination and scyther tool-based formal security verification, they have proven their scheme secured against the wide range of attacks.

In 2024, Algarni and Jan proposed a robust yet lightweight security mechanism utilizing a fuzzy extractor and the MD5 (Message Digest 5) algorithm to authenticate all IoD participants and ensure secure communication. However, the MD5 hash algorithm is widely recognized as less secure compared to more robust alternatives like the Secure Hash Algorithm (SHA-256). Consequently, the overall strength of their scheme is compromised. Moreover, their approach does not incorporate support for blockchain implementation.

Wazid et al. in 2025 developed a protocol for authentication and secure key management for internet of drone applications, the scheme is an advanced version of their protocol, with more security and efficiency.

\noindent\textit{Research gaps and novelty:} Blockchain technology offers powerful solutions to strengthen the security of the IoD environment. By enabling the creation of unique digital identities for individual drones, which are securely stored and managed on the blockchain, it helps mitigate the risk of impersonation attacks. In addition, the data coming securely from the drones to the ground station server is used for the transactions and later, the blocks formed from the authentic and genuine data from the drones are stored in the blockchain network maintained by the cloud servers. Storing data on semi-trusted cloud servers raises serious concerns about data poisoning attacks, which can significantly impact businesses and organizations by corrupting big data analytics, leading to financial losses and reputational damage. Research shows notable improvements in accuracy, recall, precision, and F1-score when data is free from poisoning attacks and is directly sourced from the blockchain. In this context, authentication among drones and other entities in the IoD environment becomes critical to ensure that genuine data is stored on the blockchain.

The literature review highlights that most existing authentication techniques for IoT and IoD environments are either inefficient in terms of communication and computation or vulnerable to various attacks. This underscores the need for a reliable and secure authenticated key agreement protocol to facilitate secure data aggregation at ground station servers in the IoD environment, with blockchain technology providing enhanced secure storage. Therefore, the objective of this work is to develop a novel and secure blockchain-based authentication and key management mechanism for IoD applications that is not only resistant to various attacks but also efficient in communication and computational costs, making it suitable for real-world practical applications.